\section{Conclusion}

This paper introduces a semantic-orientation framework for selecting, designing, and evaluating LLM bidding-agent objectives in a stylized electricity-market simulation. The experiment varies only the objective block of a modular prompt while holding the market environment, information set, response format, action space, model configuration, and TARJ structure fixed. SBERT is used to measure the semantic direction of that intervention along a profit-versus-compliance axis.

The main contribution of this paper is a framework for using prompt objectives as controlled interventions and semantic-orientation analysis as a way to evaluate their influence on LLM-agent behavior. In the simulated market, more profit-oriented prompting leads to higher markups and prices relative to the perfect competition benchmark, especially during peak and near-scarcity hours. SBERT orientation regressions support the same directional pattern. Generated TARJ reasoning orientation also moves with bidding posture. Table~\ref{tab:dose_response_summary} shows that a 0.1 increase in prompt orientation is associated with a 0.029 increase in mean TARJ reasoning orientation, with $R^2=0.967$ and $p<0.001$. This suggests that semantic analysis can help verify whether prompt meaning, agent-generated reasoning, and submitted bids are directionally aligned.

Future work can proceed in several directions. Our framework can be tested in more realistic power-market environments with richer operational constraints, firm portfolios, uncertainty, storage, renewable generation, and repeated market interaction. Future studies can also formalize the evaluation framework and test its robustness across additional LLM models, embedding methods, prompt grids, and market designs. Together, these extensions would clarify how general the observed prompt-orientation effects are and how semantic measurement can be used to systematically evaluate LLM-agent behavior in electricity-market settings.